\documentclass{article}
\usepackage{amsmath}
\usepackage{amssymb}
\usepackage[utf8]{inputenc}
\usepackage{CJKutf8}

\title{方程定义文档}
\author{Equation Compiler}
\date{\today}

\begin{document}
\begin{CJK}{UTF8}{gbsn}

\maketitle

\tableofcontents
\newpage

\section{级联计算演示 (CASCADE)}

展示方程之间的 DAG 依赖关系

\subsection{方程}

\subsubsection{stage1: 一阶变换}

\begin{equation}
    y1 = p1 \times x
\end{equation}

\subsubsection{stage2: 二阶变换}

\begin{equation}
    y2 = y1 + p2 \times x
\end{equation}

\subsubsection{stage3: 最终输出}

\begin{equation}
    y_{final} = \sqrt{|y2|}
\end{equation}

\section{数学运算演示 (MATH_DEMO)}

演示 equation-compiler 支持的各种数学运算符

\subsection{方程}

\subsubsection{damped_oscillation: 阻尼振荡}

\begin{equation}
    y_{damped} = p1 \times e^{-p3 \times t} \times \sin(2 \times \pi \times p2 \times t)
\end{equation}

\subsubsection{tanh_activation: 双曲正切激活}

\begin{equation}
    y_{tanh} = \tanh(p1 \times x)
\end{equation}

\subsubsection{angle_calc: 角度计算}

\begin{equation}
    theta = \text{atan2}(x, t)
\end{equation}

\subsubsection{rounding_demo: 取整运算}

\begin{equation}
    y_{round} = \lfloor x \rfloor + \lceil \frac{t}{2} \rceil
\end{equation}

\subsubsection{sign_abs_demo: 符号与绝对值}

\begin{equation}
    y_{sign} = \text{sgn}(x) \times \sqrt{|x|}
\end{equation}

\subsubsection{root_log_demo: 立方根与对数}

\begin{equation}
    y_{log} = \sqrt[3]{x} + \log_{2}(t + 1)
\end{equation}

\subsubsection{relu_activation: ReLU激活}

\begin{equation}
    y_{relu} = \begin{cases} x & \text{if } x > 0 \\ 0 & \text{otherwise} \end{cases}
\end{equation}

\subsubsection{modulo_demo: 取余运算}

\begin{equation}
    y_{mod} = \lfloor t \times 10 \rfloor \mod 3
\end{equation}

\subsubsection{inverse_trig: 反三角函数}

\begin{equation}
    y_{asin} = 2 \times \arcsin(\frac{x}{|x| + 1})
\end{equation}

\subsubsection{exp_growth: 指数增长}

\begin{equation}
    y_{exp} = e^{p3 \times t}
\end{equation}

\section{光合作用 (PHOTO)}

计算叶片光合速率

\subsection{方程}

\subsubsection{PHOTO-01: 动态Pmax}

\begin{equation}
    Pmax_{l} = p1 + p2 \times reserve_{ratio}
\end{equation}

\textit{参考: Lescourret 1998}

\subsubsection{PHOTO-02: Higgins光响应}

\begin{equation}
    A_{leaf} = Pmax_{l} + p3 \times 1 - e^{\frac{-p4 \times ppfd}{Pmax_{l} + p3}} - p3
\end{equation}

\textit{参考: Higgins 1992, Eq. 3}

\end{CJK}
\end{document}